\documentclass[10pt]{beamer}

\usetheme{Madrid}
\usecolortheme{whale}

\usepackage[utf8]{inputenc}
\usepackage[T1]{fontenc}
\usepackage[italian]{babel}
\usepackage{tikz}
\usetikzlibrary{arrows.meta, positioning, shapes.geometric, calc, fit, shapes.misc, shadows}
\usepackage{listings}
\usepackage{xcolor}

% --- Configurazione Colori e Link ---
\definecolor{primary}{RGB}{38, 66, 139}
\definecolor{secondary}{RGB}{200, 50, 50}
\definecolor{codebg}{RGB}{245, 245, 245}
\definecolor{keywordcolor}{RGB}{0, 0, 255}
\definecolor{stringcolor}{RGB}{163, 21, 21}
\definecolor{commentcolor}{RGB}{0, 128, 0}

\lstdefinelanguage{javascript}{
    keywords={async, await, break, case, catch, class, const, continue, debugger, default, delete, do, else, export, extends, false, finally, for, function, if, import, in, instanceof, new, null, return, super, switch, this, throw, true, try, typeof, var, void, while, with, yield, let},
    morekeywords={console, log, Error, fetch, JSON, parse, stringify, WebSocket, on, send, RTCPeerConnection, addTrack, ontrack, createOffer, setLocalDescription, setRemoteDescription},
    sensitive=true,
    morecomment=[l]{//},
    morecomment=[s]{/*}{*/},
    morestring=[b]',
    morestring=[b]",
    morestring=[b]`,
    keywordstyle=\color{primary}\bfseries,
    commentstyle=\color{commentcolor}\itshape,
    stringstyle=\color{stringcolor},
    basicstyle=\ttfamily\tiny
}

\lstset{
    basicstyle=\ttfamily\tiny,
    backgroundcolor=\color{gray!5},
    frame=single,
    rulecolor=\color{gray!30},
    keywordstyle=\color{primary}\bfseries,
    commentstyle=\color{commentcolor},
    stringstyle=\color{stringcolor},
    showstringspaces=false,
    breaklines=true
}

\setbeamercolor{title}{fg=white, bg=primary}
\setbeamercolor{frametitle}{fg=white, bg=primary}
\setbeamercolor{structure}{fg=primary}

\title[Lezione 6]{Programmazione e Tecnologie Web}
\subtitle{Real-time Communication \& Multimedia}
\author{Giuseppe Capasso}
\institute{ITS - Programmazione Web}
\date{A.A. 2025/2026}

\begin{document}

\begin{frame}
    \titlepage
\end{frame}

\begin{frame}{Sommario}
    \tableofcontents
\end{frame}

\section{Evoluzione Real-time}

\begin{frame}{Il Limite di HTTP: Request-Response}
    \begin{columns}
        \begin{column}{0.5\textwidth}
            \begin{itemize}
                \item Unidirezionale (Client $\rightarrow$ Server).
                \item Stateless (di base).
                \item Alto overhead (Header per ogni richiesta).
            \end{itemize}
        \end{column}
        \begin{column}{0.5\textwidth}
            \centering
            \begin{tikzpicture}[scale=0.7, every node/.style={scale=0.7}]
                \node[draw, thick, fill=blue!10, minimum width=1.5cm] (C) {Client};
                \node[draw, thick, fill=gray!10, minimum width=1.5cm, right=2cm of C] (S) {Server};
                \draw[->, thick, primary] ([yshift=5mm]C.east) -- node[above, font=\tiny]{Request} ([yshift=5mm]S.west);
                \draw[<-, thick, gray] ([yshift=-5mm]C.east) -- node[below, font=\tiny]{Response} ([yshift=-5mm]S.west);
            \end{tikzpicture}
        \end{column}
    \end{columns}
\end{frame}

\begin{frame}{Polling vs Long Polling}
    \centering
    \begin{tikzpicture}[
        node distance=0.5cm,
        line/.style={draw, thick, -latex},
        label/.style={font=\tiny, align=center}
    ]
        % Polling
        \node (c1) {Client};
        \node (s1) [right=2.5cm of c1] {Server};
        \draw [thick] (c1.south) -- +(0,-3cm);
        \draw [thick] (s1.south) -- +(0,-3cm);
        \node [above=0.1cm of c1, font=\tiny\bfseries] {Short Polling};
        
        \draw [line] (c1.center |- 0,-0.5) -- node[label, above] {C'è posta?} (s1.center |- 0,-0.5);
        \draw [line, gray, dashed] (s1.center |- 0,-1.0) -- node[label, below] {No.} (c1.center |- 0,-1.0);
        
        \draw [line] (c1.center |- 0,-1.7) -- node[label, above] {C'è posta?} (s1.center |- 0,-1.7);
        \draw [line, secondary] (s1.center |- 0,-2.2) -- node[label, below] {Sì, ecco i dati!} (c1.center |- 0,-2.2);

        % Long Polling
        \begin{scope}[xshift=5.5cm]
            \node (c2) {Client};
            \node (s2) [right=2.5cm of c2] {Server};
            \draw [thick] (c2.south) -- +(0,-3cm);
            \draw [thick] (s2.south) -- +(0,-3cm);
            \node [above=0.1cm of c2, font=\tiny\bfseries] {Long Polling};
            
            \draw [line] (c2.center |- 0,-0.5) -- node[label, above] {Fammi sapere quando\\ci sono novità} (s2.center |- 0,-0.5);
            \node [right=0.1cm of s2, yshift=-1.2cm, font=\tiny, primary] {Attesa...};
            \draw [line, secondary] (s2.center |- 0,-2.5) -- node[label, below] {Ecco le novità!} (c2.center |- 0,-2.5);
        \end{scope}
    \end{tikzpicture}
\end{frame}

\section{WebSocket}

\begin{frame}{WebSocket: La Svolta Bidirezionale}
    \begin{itemize}
        \item \textbf{Protocollo TCP}: Indipendente da HTTP (ma inizia via HTTP).
        \item \textbf{Full-Duplex}: Client e Server parlano contemporaneamente.
        \item \textbf{Low Overhead}: Solo 2-10 byte di frame dopo l'handshake.
    \end{itemize}
    \centering
    \begin{tikzpicture}[
        node distance=0.5cm,
        arrow/.style={-latex, thick},
        label/.style={font=\tiny}
    ]
        \node (c) {Client};
        \node (s) [right=4cm of c] {Server};
        \draw [thick] (c.south) -- +(0,-2.5cm);
        \draw [thick] (s.south) -- +(0,-2.5cm);
        
        \draw [arrow, primary] (c.center |- 0,-0.5) -- node[label, above] {HTTP Upgrade Request} (s.center |- 0,-0.5);
        \draw [arrow, gray, dashed] (s.center |- 0,-1.0) -- node[label, below] {101 Switching Protocols} (c.center |- 0,-1.0);
        
        \node[draw, primary, fill=blue!5, fit={(c.center |- 0,-1.5) (s.center |- 0,-2.2)}, inner sep=0pt, opacity=0.3] {};
        \draw [latex-latex, ultra thick, primary] (c.center |- 0,-1.8) -- node[label, above, black] {Canale WebSocket Persistente} (s.center |- 0,-1.8);
    \end{tikzpicture}
\end{frame}

\begin{frame}[fragile]{Esempio: Client \& Server (ws)}
    \begin{columns}
        \begin{column}{0.5\textwidth}
            \centering \textbf{Client (JS)}
            \begin{lstlisting}[language=javascript]
const ws = new WebSocket('ws://host');

ws.onopen = () => {
  ws.send('Ciao Server!');
};

ws.onmessage = (e) => {
  console.log('Ricevuto:', e.data);
};
            \end{lstlisting}
        \end{column}
        \begin{column}{0.5\textwidth}
            \centering \textbf{Server (Node.js)}
            \begin{lstlisting}[language=javascript]
const WebSocket = require('ws');
const wss = new WebSocket.Server({port:80});

wss.on('connection', (ws) => {
  ws.on('message', (msg) => {
    // Broadcast a tutti
    wss.clients.forEach(c => {
      c.send(msg);
    });
  });
});
            \end{lstlisting}
        \end{column}
    \end{columns}
\end{frame}

\begin{frame}{Socket.io: Più di un WebSocket}
    \begin{itemize}
        \item \textbf{Fallback}: Se WS non è disponibile, usa Polling.
        \item \textbf{Rooms}: Organizzazione logica dei client.
        \item \textbf{Auto-reconnect}: Gestione trasparente delle disconnessioni.
    \end{itemize}
    \centering
    \begin{tikzpicture}
        \node[draw, rounded corners, primary, thick, fill=blue!5, minimum width=3cm] (SIO) {Socket.io};
        \node[draw, dashed, gray, below=0.5cm of SIO, xshift=-1cm] (WS) {WebSocket};
        \node[draw, dashed, gray, below=0.5cm of SIO, xshift=1cm] (LP) {Long Polling};
        \draw[->, thick, gray] (SIO) -- (WS);
        \draw[->, thick, gray] (SIO) -- (LP);
        \node[right=0.5cm of SIO, font=\tiny, align=left] {Astrae il trasporto\\per garantire la\\connessione};
    \end{tikzpicture}
\end{frame}

\section{WebRTC}

\begin{frame}{WebRTC: Peer-to-Peer (P2P)}
    \begin{itemize}
        \item \textbf{Audio/Video/Data}: Senza passare dal server (bassa latenza).
        \item \textbf{Browser-native}: Nessun plugin richiesto.
        \item \textbf{Sicuro}: Crittografia obbligatoria (DTLS/SRTP).
    \end{itemize}
    \centering
    \begin{tikzpicture}[scale=0.8]
        \node[draw, circle, primary, inner sep=2pt] (A) {Peer A};
        \node[draw, circle, primary, inner sep=2pt, right=3cm of A] (B) {Peer B};
        \node[draw, dashed, gray, above=1.5cm of {$(A)!0.5!(B)$}] (S) {Signaling (WS)};
        
        \draw[<->, dashed, gray] (A) -- (S);
        \draw[<->, dashed, gray] (B) -- (S);
        \draw[<->, ultra thick, secondary] (A) -- node[below, font=\tiny]{Direct P2P Stream} (B);
    \end{tikzpicture}
\end{frame}

\begin{frame}{Sfida: L'attraversamento del NAT}
    \begin{itemize}
        \item \textbf{STUN}: "Ehi Server, qual è il mio IP pubblico?"
        \item \textbf{TURN}: Se il P2P fallisce, il server fa da relè (Relay).
        \item \textbf{Symmetric NAT}: La porta esterna cambia per ogni destinazione.
    \end{itemize}
    \centering
    \begin{tikzpicture}[scale=0.6, every node/.style={scale=0.6}]
        \node[draw, thick, fill=blue!5] (L) {Local Client};
        \node[draw, chamfered rectangle, fill=red!10, right=of L] (N) {NAT};
        \node[draw, thick, fill=gray!5, ellipse, right=1cm of N] (I) {Internet};
        \node[draw, thick, fill=yellow!10, above right=0.5cm of I] (STUN) {STUN Server};
        
        \draw[->] (L) -- (N);
        \draw[->, primary] (N) -- (I);
        \draw[->, dashed] (I) -- (STUN);
    \end{tikzpicture}
\end{frame}

\begin{frame}{ICE Connection States}
    \centering
    \begin{tikzpicture}[
        node distance=1.5cm,
        state/.style={draw, thick, rounded corners, fill=primary!10, minimum width=2.2cm, font=\tiny\bfseries},
        arrow/.style={-latex, thick, gray}
    ]
        \node (new) [state] {NEW};
        \node (checking) [state, right=of new] {CHECKING};
        \node (connected) [state, right=of checking] {CONNECTED};
        \node (completed) [state, right=of connected] {COMPLETED};
        \node (failed) [state, below=of checking, fill=red!10] {FAILED};

        \draw [arrow] (new) -- (checking);
        \draw [arrow] (checking) -- (connected);
        \draw [arrow] (connected) -- (completed);
        \draw [arrow] (checking) -- (failed);
    \end{tikzpicture}
    \vspace{0.3cm}
    \begin{itemize}
        \item \textbf{Checking}: Verificando i candidati (Cerca il percorso).
        \item \textbf{Connected}: Percorso trovato, dati in transito.
        \item \textbf{Failed}: NAT troppo restrittivi $\rightarrow$ Serve TURN.
    \end{itemize}
\end{frame}

\section{Esercitazione}

\begin{frame}{Architettura del Laboratorio}
    \centering
    \begin{tikzpicture}[node distance=1.5cm]
        \node[draw, thick, fill=blue!5] (UI) {Frontend (Tailwind)};
        \node[draw, thick, fill=gray!10, below=1cm of UI] (SRV) {Server (ws/Node.js)};
        \node[draw, thick, fill=green!5, left=1.5cm of SRV] (RTC1) {WebRTC Peer 1};
        \node[draw, thick, fill=green!5, right=1.5cm of SRV] (RTC2) {WebRTC Peer 2};
        
        \draw[<->, primary, thick] (UI) -- (SRV);
        \draw[<->, dashed, secondary, thick] (RTC1) to[bend left] (RTC2);
        
        \node[below=0.1cm of SRV, font=\tiny] {Simulazione locale e broadcast};
    \end{tikzpicture}
\end{frame}

\begin{frame}{Conclusioni}
    \begin{alertblock}{Takeaway}
        Il Web moderno non è più una sequenza di pagine statiche, ma un ecosistema di flussi di dati vivi.
    \end{alertblock}
    \begin{itemize}
        \item \textbf{WebSocket}: Per app centralizzate (Chat, Game server).
        \item \textbf{WebRTC}: Per app decentralizzate (Video call, P2P file sharing).
    \end{itemize}
\end{frame}

\end{document}
